\documentclass[11pt]{article}
\usepackage[margin=0.95in]{geometry}
\usepackage{amsmath,amssymb,amsthm}
\usepackage{enumitem}
\usepackage{microtype}
\usepackage{hyperref}
\hypersetup{colorlinks=true,urlcolor=blue,linkcolor=black,citecolor=black}

\setlength{\parskip}{4pt plus 1pt}
\setlength{\parindent}{0pt}

\newtheorem{thm}{Theorem}
\newtheorem{lem}[thm]{Lemma}
\newtheorem{prop}[thm]{Proposition}
\theoremstyle{definition}
\newtheorem{defn}[thm]{Definition}

\title{The Laurent-tail kill at $n=7$:\\
       seven layer-0 survivors and how they die}
\author{Jony \\ \small{(notes towards locality at $n\ge 7$)}}
\date{}

\begin{document}

\maketitle

\begin{abstract}
At $n$ external legs, the layer-0 (``Step-1'') hidden-zero kill mechanism
eliminates the coefficient $a_M$ of every chord multiset $M$ that, at
some zone $Z_{r,r+2}$, omits the eliminated vertex $r{+}1$ and avoids
the special chord $X_{r,r+2}$. Through $n=6$ this is enough: the
mechanism kills every non-triangulation. At $n=7$, however, exactly
seven non-triangulation multisets escape -- a cyclic family of ``fish.''
We show that all seven die at the next Laurent layer. For each fish
$M$ there is a zone $Z$ at which the order $X_S^{-3}$ Laurent expansion
of $1/\prod X_c$ produces a free-variable fingerprint shared by exactly
six ansatz monomials, and all five cousins of $M$ are killed by layer~0
at a single neighbouring zone. The fingerprint equation collapses to
$a_M=0$. Verification is symbolic, performed by the companion script
\verb|cascade_kill_n7.py| at \cite{repo}, with full output in
\verb|results_cascade_n7.txt|.
\end{abstract}

%--------------------------------------------------------------------------
\section{Why the fish exist: layer-0 misses seven heptagon multisets}
%--------------------------------------------------------------------------

We work with a heptagon ($n=7$) with vertices $1,2,\dots,7$ in cyclic
order. Chords are pairs $(i,j)$ with $1\le i<j\le 7$ that are not
polygon edges. The general planar rational ansatz of mass dimension
$-2(n{-}3)=-8$ is
\[
   B \;=\; \sum_{M} \frac{a_M}{\prod_{c\in M} X_c},
\]
where $M$ ranges over multisets of $N := n{-}3 = 4$ chords. A multiset
is a \emph{triangulation} iff its four chords are pairwise non-crossing;
the $\mathrm{Tr}(\phi^3)$ tree amplitude is a sum over the $C_5=42$
heptagon triangulations.

\paragraph{Zones and layer-0 kill.}
For each $r\in\{1,\dots,n\}$, the zone $Z_{r,r+2}$ is the codimension-2
hidden-zero locus on which the chords incident to vertex $r{+}1$ become
linear combinations of chords incident to $r$ and the special chord
$S=X_{r,r+2}$. Explicitly,
\[
   X_{r+1,k} \;=\; X_{r,k} - S \qquad (k\ne r,r{+}1,r{+}2),
\]
and the bare relation $X_{r+1,r-1} = -S$. At zone $Z_{r,r+2}$, every
chord is exactly one of: special ($S$), bare ($X_{r+1,r-1}$),
non-bare substituted ($X_{r+1,k}$, with companion $X_{r,k}$ free), or
plain free (every chord avoiding both $r$ and $r{+}1$). Sending
$S\to\infty$ and reading the leading Laurent order, linear independence
of the free-variable monomials forces $a_M=0$ whenever
\begin{itemize}[leftmargin=*,topsep=2pt,parsep=0pt,itemsep=2pt]
\item[(K1)] $M$ contains neither $S$ nor the bare;
\item[(K2)] for every (companion, substitute) pair, $M$ does not contain
       both members.
\end{itemize}
We write $M\in\mathcal{K}_{r,r+2}$ when (K1)+(K2) hold.

\paragraph{The seven layer-0 survivors at $n=7$.}
A complete enumeration\footnote{Performed by
\texttt{step1\_uncaught.py} in the companion repository \cite{repo}; the
output is \texttt{nonlocal\_n7\_step1\_survivors.pdf}.} shows that
exactly seven non-triangulation multisets satisfy
$M\notin\mathcal{K}_{r,r+2}$ for \emph{every} $r$:
\begin{equation}\label{eq:fish}
\begin{aligned}
M_1 &= \{(1,3),(1,6),(2,4),(4,6)\}, \quad &
M_2 &= \{(1,3),(1,6),(3,5),(4,6)\}, \\
M_3 &= \{(1,3),(1,6),(3,5),(5,7)\}, &
M_4 &= \{(1,3),(2,7),(3,5),(5,7)\}, \\
M_5 &= \{(1,6),(2,4),(2,7),(4,6)\}, &
M_6 &= \{(2,4),(2,7),(3,5),(5,7)\}, \\
M_7 &= \{(2,4),(2,7),(4,6),(5,7)\}. & &
\end{aligned}
\end{equation}
They are the seven cyclic rotations of $M_1$ under $v\mapsto v{+}\sigma$
mod $7$ for $\sigma=0,2,4,6,5,1,3$ respectively. Visually each is a
``fish'' of two crossing chords plus a containing pair of chords (cf.\
the survivor diagram in \cite{repo}).

\begin{prop}[Layer-0 fails on every fish]\label{prop:layer0fails}
For every $i\in\{1,\dots,7\}$ and every $r\in\{1,\dots,7\}$,
$M_i\notin\mathcal{K}_{r,r+2}$.
\end{prop}

\begin{proof}
By cyclic symmetry it suffices to do $i=1$. The vertices used by $M_1$
are $\{1,2,3,4,6\}$; vertices $5$ and $7$ are unused. For
$M_1\in\mathcal{K}_{r,r+2}$, condition (K1) requires the special
$X_{r,r+2}\notin M_1$ and the bare $X_{r+1,r-1}\notin M_1$. Test each
$r$:
\begin{itemize}[leftmargin=*,topsep=1pt,parsep=0pt,itemsep=1pt]
\item $r=1$: special $X_{1,3}\in M_1$; (K1) fails.
\item $r=2$: special $X_{2,4}\in M_1$; (K1) fails.
\item $r=3$: special $X_{3,5}\notin M_1$, bare $X_{4,2}=(2,4)\in M_1$;
       (K1) fails.
\item $r=4$: special $X_{4,6}=(4,6)\in M_1$; (K1) fails.
\item $r=5$: special $X_{5,7}\notin M_1$, bare $X_{6,4}=(4,6)\in M_1$;
       (K1) fails.
\item $r=6$: special $X_{6,1}=(1,6)\in M_1$; (K1) fails.
\item $r=7$: special $X_{7,2}=(2,7)\notin M_1$, bare
       $X_{1,6}=(1,6)\in M_1$; (K1) fails.
\end{itemize}
So (K1) fails at every zone, and $M_1\notin\mathcal{K}_{r,r+2}$ for all
$r$. \qed
\end{proof}

The proof shows the structural reason: each fish either uses the special
or contains the bare at every cyclic zone. This is what makes the fish
``unbreakable'' by layer-0 -- they are precisely the multisets whose
chords cyclically saturate the special/bare slots.

%--------------------------------------------------------------------------
\section{The Laurent-tail kill}
%--------------------------------------------------------------------------

\paragraph{The mechanism.}
At zone $Z_{r,r+2}$ with $S=X_{r,r+2}$, the geometric-series identity
\begin{equation}\label{eq:geom}
   \frac{1}{X_{r,k}-S} \;=\; -\frac{1}{S} \;-\; \frac{X_{r,k}}{S^2}
   \;-\; \frac{X_{r,k}^{2}}{S^3} \;-\; \cdots
\end{equation}
expands every non-bare substitute $1/X_{r+1,k}$ as a Laurent series in
$1/S$. The bare $1/X_{r+1,r-1}=-1/S$ has no tail. Multiplying $N$ such
factors gives, for each ansatz monomial, a Laurent series in $1/S$ whose
coefficients are rational \emph{free-variable} monomials -- monomials
in the companions $\{X_{r,k}\}$ and the born-free chords (those avoiding
both $r$ and $r{+}1$). Algebraic independence of the free variables
implies that every free-variable monomial at every order in $S$ has
total coefficient zero on $B|_{Z_{r,r+2}}=0$.

\paragraph{Notation.}
A \emph{fingerprint} is a specific rational free-variable monomial
$\mathcal F$. The \emph{fingerprint equation} of $(\mathcal F, k, r)$ is
the equation that the order-$S^{-(N+k)}$ coefficient of $\mathcal F$ in
$B|_{Z_{r,r+2}}$ equals zero. Layer~$k$ refers to subleading order $k$
beyond layer~0.

\paragraph{The kill of $M_1$.}
Take $M_1=\{(1,3),(1,6),(2,4),(4,6)\}$ at zone $Z_{5,7}$, where
$S=X_{5,7}$, eliminated vertex $6$, and the substitutions are
\begin{equation}\label{eq:Z57subs}
   X_{1,6}=X_{1,5}-S,\quad X_{2,6}=X_{2,5}-S,\quad
   X_{3,6}=X_{3,5}-S,\quad X_{4,6}=-S.
\end{equation}
The chords $(1,3)$ and $(2,4)$ avoid both $5$ and $6$, hence are free.
The Laurent expansion of the fish monomial on $Z_{5,7}$ is
\[
   \frac{1}{X_{1,3}\,X_{2,4}\,X_{1,6}\,X_{4,6}}
   \;=\;
   \frac{1}{X_{1,3}X_{2,4}}\cdot\frac{1}{X_{1,5}-S}\cdot\frac{-1}{S}.
\]
Using \eqref{eq:geom}, the order $S^{-3}$ piece is
$\dfrac{X_{1,5}}{X_{1,3}X_{2,4}}\cdot\dfrac{1}{S^{3}}$. We isolate the
fingerprint
\[
   \mathcal F_1 \;:=\; \frac{X_{1,5}}{X_{1,3}X_{2,4}}.
\]

\begin{lem}[Six-monomial fingerprint equation for $M_1$]\label{lem:eqM1}
The fingerprint $\mathcal F_1$ at order $S^{-3}$ on $Z_{5,7}$ receives
contributions from exactly six ansatz monomials. The fingerprint
equation is
\begin{equation}\label{eq:eqM1}
\hspace*{-2pt}
\begin{aligned}
0 \;=\; & -\,a_{\{(1,3),(1,5),(1,6),(2,4)\}}
        \;+\;2\,a_{\{(1,3),(1,6),(1,6),(2,4)\}}
        \;+\;a_{\{(1,3),(1,6),(2,4),(2,6)\}} \\
        & +\,a_{\{(1,3),(1,6),(2,4),(3,6)\}}
        \;+\;a_{\{(1,3),(1,6),(2,4),(4,6)\}}
        \;-\;a_{\{(1,3),(1,6),(2,4),(5,7)\}}.
\end{aligned}
\end{equation}
\end{lem}

\begin{proof}
A monomial of $B$ contributes a power of $X_{1,5}$ in the numerator at
order $S^{-3}$ iff one of its chord factors is incident to vertex $6$
on the substituted side -- specifically iff it contains $X_{1,6}$,
because $X_{1,5}$ enters only through the Laurent tail
\eqref{eq:Z57subs} of $1/X_{1,6}$. (The companions of $X_{2,6}$ and
$X_{3,6}$ supply $X_{2,5}$ and $X_{3,5}$ respectively, not $X_{1,5}$.)
Likewise the denominator of $\mathcal F_1$ requires the monomial to
contain $X_{1,3}$ and $X_{2,4}$ each with multiplicity $1$. A size-$4$
multiset thus has chord factors $\{X_{1,3}, X_{2,4}, X_{1,6}, X_a\}$ for
some fourth chord $X_a$. Direct expansion using \eqref{eq:geom} gives
each of the six contributing $X_a\in\{X_{1,5}, X_{1,6}, X_{2,6}, X_{3,6},
X_{4,6}, X_{5,7}\}$ the integer scalar shown. The full computation is
performed by \texttt{fingerprint\_equation} in
\verb|cascade_kill_n7.py|. \qed
\end{proof}

\begin{lem}[All five cousins die by layer-0 at $Z_{4,6}$]\label{lem:cousins}
Each of the five non-fish multisets in \eqref{eq:eqM1} satisfies
$M\in\mathcal{K}_{4,6}$, hence its coefficient is killed by layer-0 at
zone $Z_{4,6}$.
\end{lem}

\begin{proof}
At $Z_{4,6}$: special $X_{4,6}=(4,6)$, bare $X_{5,3}=(3,5)$,
companion-substitute pairs $(X_{4,7},X_{5,7})$, $(X_{4,1},X_{5,1})$
$=(X_{1,4},X_{1,5})$, $(X_{4,2},X_{5,2})=(X_{2,4},X_{2,5})$. We check
(K1)+(K2) for each cousin:
\begin{itemize}[leftmargin=*,topsep=1pt,parsep=0pt,itemsep=1pt]
\item $\{(1,3),(1,5),(1,6),(2,4)\}$: special $\notin$, bare $\notin$,
       pair $(X_{1,4},X_{1,5})$ has only $X_{1,5}$, others have at most
       one. \checkmark
\item $\{(1,3),(1,6),(1,6),(2,4)\}$: special $\notin$, bare $\notin$,
       no pair has both members. \checkmark
\item $\{(1,3),(1,6),(2,4),(2,6)\}$: special $\notin$, bare $\notin$,
       no pair has both ($X_{2,4}\in$, $X_{2,5}\notin$).
       \checkmark
\item $\{(1,3),(1,6),(2,4),(3,6)\}$: special $\notin$, bare $\notin$,
       no pair has both. \checkmark
\item $\{(1,3),(1,6),(2,4),(5,7)\}$: special $\notin$, bare $\notin$,
       pair $(X_{4,7},X_{5,7})$ has only $X_{5,7}$, pair
       $(X_{2,4},X_{2,5})$ has only $X_{2,4}$. \checkmark
\end{itemize}
All five satisfy (K1)+(K2) at $Z_{4,6}$. \qed
\end{proof}

\begin{thm}[Cascade kill of all seven $n=7$ fish]\label{thm:main}
For each $i\in\{1,\dots,7\}$, $a_{M_i}=0$.
\end{thm}

\begin{proof}
For $i=1$, substitute Lemma~\ref{lem:cousins} into \eqref{eq:eqM1}:
all five cousin coefficients are zero, so the equation collapses to
$+a_{M_1}=0$. For $i=2,\dots,7$, apply the cyclic shift
$v\mapsto v{+}\sigma_i$ mod $7$ to the entire $i=1$ argument: the recipe
$(Z_{5,7},\mathcal F_1)$ becomes $(Z_{5+\sigma_i,7+\sigma_i},
\mathrm{rot}_{\sigma_i}(\mathcal F_1))$, the six-monomial equation
becomes the cyclic rotate of \eqref{eq:eqM1}, and the cousins'
$Z_{4,6}$-layer-0 kill becomes $Z_{4+\sigma_i,6+\sigma_i}$-layer-0 kill.
The companion script \verb|cascade_kill_n7.py| recomputes the rotated
equation directly (rather than rotating the symbolic answer), confirming
the recipe and verifying every cousin kill at the rotated zone; see
\verb|results_cascade_n7.txt|. \qed
\end{proof}

%--------------------------------------------------------------------------
\section{Companion code and reproducibility}\label{sec:code}
%--------------------------------------------------------------------------

The script
\href{https://github.com/Jony-GydeRnE/locality-proof/blob/main/cascade_kill_n7.py}%
{\tt cascade\_kill\_n7.py} reproduces every step of
Theorem~\ref{thm:main} symbolically. The main routines:

\begin{itemize}[leftmargin=*,topsep=2pt,parsep=0pt,itemsep=2pt]
\item \texttt{laurent\_coefficient(M, r, n, k)}: Laurent expansion of
       $1/\prod_{c\in M} X_c$ on $Z_{r,r+2}$ to order $S^{-k}$, via
       \texttt{sympy.series} after the substitution $S\mapsto 1/u$.
\item \texttt{coefficient\_of\_monomial(expr, target, free\_vars)}:
       extract the scalar coefficient of a target free-variable
       monomial from a sum of rational monomials.
\item \texttt{candidate\_multisets\_for\_fingerprint(fp, n, N)}: fast
       pre-filter that enumerates only multisets whose chord factors are
       compatible with producing \texttt{fp} (so the search runs in
       seconds rather than hours).
\item \texttt{fingerprint\_equation(M, r, n, k, fp, pool)}: returns
       $[(M', \text{scalar})]$ for every multiset that contributes the
       fingerprint at order $S^{-k}$ on $Z_{r,r+2}$.
\item \texttt{verify\_each\_fish\_from\_scratch}: for each fish,
       computes its kill recipe by cyclic rotation of the $M_1$ recipe
       and \emph{verifies it directly} by recomputing the fingerprint
       equation and checking that every cousin is layer-0 killable.
\item \texttt{verify\_all\_seven\_fish}: top-level driver, writes the
       human-readable cascade trace to \texttt{results\_cascade\_n7.txt}.
\end{itemize}

The output file \texttt{results\_cascade\_n7.txt} contains, for each
fish, the kill zone and order, the fingerprint, the six-term equation,
and the layer-0 kill zone of every cousin. All seven kills are
confirmed.

%--------------------------------------------------------------------------
\section{Outlook}
%--------------------------------------------------------------------------

The cascade has a strikingly simple shape: \emph{one Laurent step deeper
at one zone, with prerequisites all of layer-0 type at one neighbouring
zone}. The cyclic structure of \eqref{eq:fish} reflects this -- each
fish is the rotated image of $M_1$, and the kill recipe rotates with
it. This raises the natural conjecture that every layer-0 survivor at
$n\ge 7$ dies via a depth-1 cascade of the same shape: identify the
per-fish recipe combinatorially, then verify the fingerprint equation's
cousins are all layer-0 killable. The next experiment is to check this
at $n=8,9$ using the same script with the obvious adaptations. If the
depth-1 phenomenon persists, the path to a clean all-$n$ locality proof
is to reduce the all-layer fingerprint problem to a single-step
combinatorial induction.

\begin{thebibliography}{9}
\bibitem{repo}
GitHub repository \emph{locality-proof}:
\url{https://github.com/Jony-GydeRnE/locality-proof}.
Files: \texttt{cascade\_kill\_n7.py} (verifier),
\texttt{results\_cascade\_n7.txt} (cascade trace),
\texttt{step1\_uncaught.py} (layer-0 enumeration),
\texttt{nonlocal\_n7\_step1\_survivors.pdf} (the seven fish drawn).

\bibitem{rodina}
A.\ Rodina, \emph{Hidden zeros are equivalent to enhanced ultraviolet
scaling and lead to unique amplitudes in $\mathrm{Tr}(\phi^3)$ theory},
arXiv:2406.04234.

\bibitem{aharkanihamed}
N.\ Arkani-Hamed, Q.\ Cao, J.\ Dong, C.\ Figueiredo, S.\ He,
\emph{Hidden zeros for particle/string amplitudes and the unity of
colored scalars, pions, and gluons}, arXiv:2312.16282.
\end{thebibliography}

\end{document}
